% GENERATED FILE. Edit canonical lessons, then run npm run generate:books.
% canonical-source-hash: fnv1a64:eeafd46f33f2c3e1
% canonical-lessons: HI-C283-agla, HI-C283-ektikat, HI-C283-khona, HI-C283-chutgayi, HI-R283-at-the-station

\chapter{Tickets, Trains and Trouble}
\label{ch:hi-tickets-trains-and-trouble}
\hlchaptermodality{283}

\begin{chapteropening}
\textbf{By the end of this chapter:} \emph{I can buy a ticket and ask when the next train leaves (aglī ṭren kitne baje calegī?), and report a lost bag or a missed train (merā thailā kho gayā, merī ṭren chūṭ gaī).}

Complete HI-R283-at-the-station: buy a ticket, ask about the next train, and report what went wrong.
\end{chapteropening}

\section[\texorpdfstring{\hi{अगला} (aglā)}{अगला (aglā)}]{\hi{अगला} (aglā) — next}

\label{lesson:HI-C283-agla}



\begin{quote}
Before the new one: say \textquotedblleft{}last week\textquotedblright{}. (\textbf{\hi{पिछले} \hi{हफ़्ते}}.) Then \textbf{\hi{आगे}}, \textquotedblleft{}ahead\textquotedblright{}. This lesson's word is the partner of the first.
\end{quote}

\subsection*{You'll want to know: \hi{अगला}}
\textbf{\hi{अगला}} — \emph{aglā} — \textquotedblleft{}next\textquotedblright{}.

\begin{quote}
\textbf{\hi{अगली} \hi{ट्रेन} \hi{कब} \hi{आएगी}?}
\end{quote}

\begin{quote}
\emph{aglī ṭren kab āegī?} — \textquotedblleft{}When will the next train come?\textquotedblright{}
\end{quote}

It agrees as \textbf{\hi{पिछला}} does: \textbf{\hi{अगला} \hi{स्टेशन}} (\emph{aglā sṭeśan}), \textquotedblleft{}the next station\textquotedblright{}, but \textbf{\hi{अगली} \hi{ट्रेन}}, because \textbf{\hi{ट्रेन}} is feminine. As a time it takes the form of \textbf{\hi{पिछले} \hi{हफ़्ते}}: \textbf{\hi{अगले} \hi{हफ़्ते}} (\emph{agle hafte}), \textquotedblleft{}next week\textquotedblright{}; \textbf{\hi{अगले} \hi{साल}} (\emph{agle sāl}), \textquotedblleft{}next year\textquotedblright{}.

\begin{quote}
\textbf{\hi{अगले} \hi{हफ़्ते} \hi{हम} \hi{गाँव} \hi{जाएँगे।}}
\end{quote}

\begin{quote}
\emph{agle hafte ham gā̃v jāeṅge.} — \textquotedblleft{}Next week we will go to the village.\textquotedblright{}
\end{quote}

\begin{cousinweb}
\textbf{\hi{अगला}} (\emph{aglā}) has no row of its own in CDIAL. It is usually connected with \textbf{\hi{आगे}} (\emph{āge}), \textquotedblleft{}ahead\textquotedblright{}, which CDIAL traces to Sanskrit \emph{agra}, \textquotedblleft{}top, front\textquotedblright{} (CDIAL 68); under the same entry CDIAL lists a Hindi \emph{āgil}, \textquotedblleft{}foremost\textquotedblright{}. \textbf{\hi{पिछला}} and \textbf{\hi{अगला}} make a pair: the one behind you and the one in front.
\end{cousinweb}

\subsection*{Your turn}
Say these aloud:

\begin{itemize}
\raggedright
  \item \emph{aglā}, then \emph{aglī ṭren}
  \item \textquotedblleft{}When will the next train come?\textquotedblright{} — \emph{aglī ṭren kab āegī?}
  \item \textquotedblleft{}Last week I was at home; next week I will go to the village\textquotedblright{} — \emph{pichle hafte maiṁ ghar par thā; agle hafte maiṁ gā̃v jāūṅgā}
\end{itemize}

\begin{tcolorbox}[breakable,colback=teal!4,colframe=teal!35!black,title={Before you move on}]
What does \textbf{\hi{अगला}} mean? (\textquotedblleft{}Next\textquotedblright{}.) \textquotedblleft{}Next week\textquotedblright{}? (\textbf{\hi{अगले} \hi{हफ़्ते}}.) \textquotedblleft{}The next train\textquotedblright{}? (\textbf{\hi{अगली} \hi{ट्रेन}}: \textbf{\hi{ट्रेन}} is feminine.)
\end{tcolorbox}
\section[\texorpdfstring{\hi{शहर} \hi{के} \hi{लिए} \hi{दो} … (śahar ke liye do …)}{शहर के लिए दो … (śahar ke liye do …)}]{\hi{शहर} \hi{के} \hi{लिए} \hi{दो} \hi{टिकट} \hi{दीजिए} (śahar ke liye do ṭikaṭ dījie) — at the ticket window}

\label{lesson:HI-C283-ektikat}



\begin{quote}
You have asked for a ticket since you first learned the word: \textbf{\hi{एक} \hi{टिकट} \hi{दीजिए}}. Say it, then \textquotedblleft{}the next train\textquotedblright{}. (\emph{aglī ṭren}.) This lesson buys the right ticket for the right train.
\end{quote}

\subsection*{You'll want to know: buying a ticket}
\begin{quote}
\textbf{\hi{शहर} \hi{के} \hi{लिए} \hi{दो} \hi{टिकट} \hi{दीजिए।}}
\end{quote}

\begin{quote}
\emph{śahar ke liye do ṭikaṭ dījie.} — \textquotedblleft{}Two tickets for the city, please.\textquotedblright{}
\end{quote}

The place you are going takes \textbf{\hi{के} \hi{लिए}}, \textquotedblleft{}for\textquotedblright{}. \textbf{\hi{टिकट}} is masculine in standard Hindi, though you will hear it made feminine too, and it does not change for two.

\begin{itemize}
\raggedright
  \item \textbf{\hi{एक} \hi{टिकट} \hi{कितने} \hi{का} \hi{है}?} (\emph{ek ṭikaṭ kitne kā hai?}) — How much is one ticket?
  \item \textbf{\hi{सौ} \hi{रुपये} \hi{का} \hi{है।}} (\emph{sau rupaye kā hai.}) — It's a hundred rupees.
\end{itemize}

In \textbf{\hi{कितने} \hi{का}}, \textquotedblleft{}of how much\textquotedblright{}, the \textbf{\hi{का}} agrees with the ticket.

\subsection*{You'll want to know: when it leaves, and how long}
\begin{quote}
\textbf{\hi{अगली} \hi{ट्रेन} \hi{कितने} \hi{बजे} \hi{चलेगी}?}
\end{quote}

\begin{quote}
\emph{aglī ṭren kitne baje calegī?} — \textquotedblleft{}What time does the next train leave?\textquotedblright{}
\end{quote}

A train \textbf{\hi{चलती} \hi{है}}, \textquotedblleft{}moves off\textquotedblright{}, and \textbf{\hi{चलेगी}} is its future. \textquotedblleft{}At what time\textquotedblright{} is \textbf{\hi{कितने} \hi{बजे}}, and the answer is a time you know: \textbf{\hi{पाँच} \hi{बजे}}.

How long will it take? \textbf{\hi{लगना}} again, with the time as the thing:

\begin{itemize}
\raggedright
  \item \textbf{\hi{कितना} \hi{समय} \hi{लगेगा}?} (\emph{kitnā samay lagegā?}) — How long will it take?
  \item \textbf{\hi{दो} \hi{घंटे} \hi{लगेंगे।}} (\emph{do ghaṇṭe lageṅge.}) — It will take two hours.
\end{itemize}

\textbf{\hi{लगेगा}} follows \textbf{\hi{समय}}; \textbf{\hi{लगेंगे}} follows \textbf{\hi{दो} \hi{घंटे}}.

\subsection*{Your turn}
\begin{itemize}
\raggedright
  \item \emph{Say it:} \textquotedblleft{}It will take two hours\textquotedblright{} — \emph{do ghaṇṭe lageṅge}
\end{itemize}

\begin{tcolorbox}[breakable,colback=teal!4,colframe=teal!35!black,title={Before you move on}]
\textquotedblleft{}Two tickets for the city\textquotedblright{}? (\textbf{\hi{शहर} \hi{के} \hi{लिए} \hi{दो} \hi{टिकट} \hi{दीजिए}}.) \textquotedblleft{}How much is one ticket?\textquotedblright{} (\textbf{\hi{एक} \hi{टिकट} \hi{कितने} \hi{का} \hi{है}?}) \textquotedblleft{}What time does it leave?\textquotedblright{} (\textbf{\hi{कितने} \hi{बजे} \hi{चलेगी}?}) \textquotedblleft{}How long will it take?\textquotedblright{} (\textbf{\hi{कितना} \hi{समय} \hi{लगेगा}?})
\end{tcolorbox}
\section[\texorpdfstring{\hi{खोना} (khonā)}{खोना (khonā)}]{\hi{खोना} (khonā) — to lose, to be lost}

\label{lesson:HI-C283-khona}



\begin{quote}
Before the new one: say a bag and a key. (\emph{thailā}, \emph{cābī}.) Then \textquotedblleft{}he fell asleep\textquotedblright{}. (\emph{vah so gayā}.) The new verb works like that last one.
\end{quote}

\subsection*{You'll want to know: \hi{खोना}}
\textbf{\hi{खोना}} — \emph{khonā} — \textquotedblleft{}to lose; to be lost\textquotedblright{}.

\begin{quote}
\textbf{\hi{मेरा} \hi{थैला} \hi{खो} \hi{गया।}}
\end{quote}

\begin{quote}
\emph{merā thailā kho gayā.} — \textquotedblleft{}My bag got lost.\textquotedblright{}
\end{quote}

With \textbf{\hi{गया}} after it, as in \textbf{\hi{सो} \hi{गया}}, the lost thing is the subject, and the verb follows it: \textbf{\hi{मेरी} \hi{चाबी} \hi{खो} \hi{गई}} (\emph{merī cābī kho gaī}), \textquotedblleft{}my key got lost\textquotedblright{}, because \textbf{\hi{चाबी}} is feminine.

With \textbf{\hi{दिया}} after it, you are the one who lost it, and \textbf{\hi{ने}} comes in:

\begin{quote}
\textbf{\hi{मैंने} \hi{अपनी} \hi{चाबी} \hi{खो} \hi{दी।}}
\end{quote}

\begin{quote}
\emph{maiṁne apnī cābī kho dī.} — \textquotedblleft{}I lost my key.\textquotedblright{}
\end{quote}

\textbf{\hi{दी}} follows \textbf{\hi{चाबी}}. \textbf{\hi{अपनी}} is \textquotedblleft{}my own\textquotedblright{}, the \textbf{\hi{अपना}} you use wherever the owner is also the doer.

\begin{cousinweb}
\textbf{\hi{खोना}} (\emph{khonā}) is a tadbhava. CDIAL traces it to Sanskrit \emph{kṣapayati}, \textquotedblleft{}destroys; throws\textquotedblright{} (CDIAL 3651), and gives \textquotedblleft{}to lose, waste, throw away\textquotedblright{} for Hindi.
\end{cousinweb}

\subsection*{Your turn}
Say these aloud:

\begin{itemize}
\raggedright
  \item \emph{khonā}, then \emph{kho gayā}
  \item \textquotedblleft{}My key got lost\textquotedblright{} — \emph{merī cābī kho gaī}
  \item \textquotedblleft{}I lost my key\textquotedblright{} — \emph{maiṁne apnī cābī kho dī}
\end{itemize}

\begin{tcolorbox}[breakable,colback=teal!4,colframe=teal!35!black,title={Before you move on}]
What does \textbf{\hi{खोना}} mean? (\textquotedblleft{}To lose; to be lost\textquotedblright{}.) \textquotedblleft{}My bag got lost\textquotedblright{}? (\textbf{\hi{मेरा} \hi{थैला} \hi{खो} \hi{गया}}.) Why \textbf{\hi{खो} \hi{गई}} for a key? (\textbf{\hi{चाबी}} is feminine.)
\end{tcolorbox}
\section[\texorpdfstring{\hi{मेरी} \hi{ट्रेन} \hi{छूट} \hi{गई} (merī ṭren chūṭ gaī)}{मेरी ट्रेन छूट गई (merī ṭren chūṭ gaī)}]{\hi{मेरी} \hi{ट्रेन} \hi{छूट} \hi{गई} (merī ṭren chūṭ gaī) — when something goes wrong}

\label{lesson:HI-C283-chutgayi}



\begin{quote}
Say \textquotedblleft{}my bag got lost\textquotedblright{}. (\emph{merā thailā kho gayā}.) Then \textbf{\hi{छूटना}}, \textquotedblleft{}to be left behind\textquotedblright{}. (\emph{chūṭnā}.)
\end{quote}

\subsection*{You'll want to know: reporting a loss}
\begin{quote}
\textbf{\hi{मेरा} \hi{थैला} \hi{खो} \hi{गया।} \hi{उसमें} \hi{मेरा} \hi{बटुआ} \hi{था।}}
\end{quote}

\begin{quote}
\emph{merā thailā kho gayā. usmeṁ merā baṭuā thā.} — \textquotedblleft{}My bag got lost. My wallet was in it.\textquotedblright{}
\end{quote}

\textbf{\hi{उसमें}} is \textbf{\hi{उस}} with \textbf{\hi{में}}: \textquotedblleft{}in it\textquotedblright{}. Then ask for help, and say whom you told:

\begin{itemize}
\raggedright
  \item \textbf{\hi{मेरी} \hi{मदद} \hi{कीजिए।}} (\emph{merī madad kījiye.}) — Please help me. \textbf{\hi{मदद}} is feminine: \textbf{\hi{मेरी}}.
  \item \textbf{\hi{मैंने} \hi{पुलिस} \hi{को} \hi{बताया।}} (\emph{maiṁne pulis ko batāyā.}) — I told the police.
  \item \textbf{\hi{थाना} \hi{कहाँ} \hi{है}?} (\emph{thānā kahā̃ hai?}) — Where is the police station?
\end{itemize}

\subsection*{You'll want to know: missing a train}
\begin{quote}
\textbf{\hi{मेरी} \hi{ट्रेन} \hi{छूट} \hi{गई।}}
\end{quote}

\begin{quote}
\emph{merī ṭren chūṭ gaī.} — \textquotedblleft{}I missed my train.\textquotedblright{}
\end{quote}

Word for word, \textquotedblleft{}my train got left behind\textquotedblright{}. The train is the subject, so \textbf{\hi{छूट} \hi{गई}} follows it, and there is no \textbf{\hi{मैंने}}. With a bus, the same: \textbf{\hi{मेरी} \hi{बस} \hi{छूट} \hi{गई}} (\emph{merī bas chūṭ gaī}).

The cause goes first, with \textbf{\hi{इसलिए}}:

\begin{quote}
\textbf{\hi{मुझे} \hi{देर} \hi{हो} \hi{गई}, \hi{इसलिए} \hi{मेरी} \hi{ट्रेन} \hi{छूट} \hi{गई।}}
\end{quote}

\begin{quote}
\emph{mujhe der ho gaī, isliye merī ṭren chūṭ gaī.} — \textquotedblleft{}I was late, so I missed my train.\textquotedblright{}
\end{quote}

And what it cost you, with \textbf{\hi{पड़ना}}: \textbf{\hi{मुझे} \hi{बस} \hi{से} \hi{जाना} \hi{पड़ा}} (\emph{mujhe bas se jānā paṛā}), \textquotedblleft{}I had to go by bus\textquotedblright{}.

\subsection*{Your turn}
Say these aloud:

\begin{itemize}
\raggedright
  \item \textquotedblleft{}Please help me\textquotedblright{} — \emph{merī madad kījiye}
  \item \textquotedblleft{}I told the police\textquotedblright{} — \emph{maiṁne pulis ko batāyā}
  \item \textquotedblleft{}I was late, so I missed my bus\textquotedblright{} — \emph{mujhe der ho gaī, isliye merī bas chūṭ gaī}
\end{itemize}

\begin{tcolorbox}[breakable,colback=teal!4,colframe=teal!35!black,title={Before you move on}]
\textquotedblleft{}My bag got lost\textquotedblright{}? (\textbf{\hi{मेरा} \hi{थैला} \hi{खो} \hi{गया}}.) \textquotedblleft{}My wallet was in it\textquotedblright{}? (\textbf{\hi{उसमें} \hi{मेरा} \hi{बटुआ} \hi{था}}.) \textquotedblleft{}I missed my train\textquotedblright{}? (\textbf{\hi{मेरी} \hi{ट्रेन} \hi{छूट} \hi{गई}}.) Why no \textbf{\hi{मैंने}}? (The train is the subject: it got left behind.)
\end{tcolorbox}
\section[(review)]{Review — at the station}

\label{lesson:HI-R283-at-the-station}



\begin{quote}
Say \textquotedblleft{}the next train\textquotedblright{} and \textquotedblleft{}next week\textquotedblright{}. (\emph{aglī ṭren}, \emph{agle hafte}.) Then \textquotedblleft{}my bag got lost\textquotedblright{}. (\emph{merā thailā kho gayā}.)
\end{quote}

\subsection*{Your turn}
Say these aloud:

\begin{itemize}
\raggedright
  \item Excuse me, where is the ticket office? Go straight on, then turn left.
  \item One ticket for the village, please. How much is it?
  \item When will the next train come? How long will it take?
  \item My key got lost. Then the other way: I lost my key.
  \item I was late, so I missed my train. Please help me.
  \item Next week I will go by bus.
\end{itemize}

\begin{tcolorbox}[breakable,colback=teal!4,colframe=teal!35!black,title={Before you move on}]
\textquotedblleft{}The next train\textquotedblright{}? (\textbf{\hi{अगली} \hi{ट्रेन}}.) \textquotedblleft{}How much is one ticket?\textquotedblright{} (\textbf{\hi{एक} \hi{टिकट} \hi{कितने} \hi{का} \hi{है}?}) \textquotedblleft{}What time does it leave?\textquotedblright{} (\textbf{\hi{कितने} \hi{बजे} \hi{चलेगी}?}) \textquotedblleft{}I lost my key\textquotedblright{}? (\textbf{\hi{मैंने} \hi{अपनी} \hi{चाबी} \hi{खो} \hi{दी}}: \textbf{\hi{ने}}, and \textbf{\hi{दी}} for the key.) \textquotedblleft{}I missed my train\textquotedblright{}? (\textbf{\hi{मेरी} \hi{ट्रेन} \hi{छूट} \hi{गई}}.)
\end{tcolorbox}
